\section{The Fermions} \label{sec:fermions}

We now open one generation and read off its particles. The construction is due to Furey \cite{furey2018}, building on
G{\"u}naydin and Gürsey \cite{gunaydin1973}, and we carry it out by explicit computation on the octonion multiplication
table. The output is one full generation of quarks and leptons, with the correct color structure and the exact
fractional electric charges, and nothing else.

\subsection{Octonion multiplications make a Clifford algebra}

Although the octonions are not associative, the linear maps they induce are. For each unit $e_a$ let $L_a$ be the
operation of multiplying on the left by $e_a$, a real $8 \times 8$ matrix. We compute these seven matrices directly
from the multiplication table and find two facts.
\begin{equation}
  L_a^2 = -I, \qquad L_a L_b + L_b L_a = 0 \ \ (a \neq b), \qquad a, b \in \{1, \dots, 7\}.
\end{equation}
Each squares to minus the identity, and distinct ones anticommute. These are exactly the defining relations of the
Clifford algebra $C\ell(0,7)$, the algebra of seven anticommuting square roots of $-1$. So the octonion left-multiplications
generate $C\ell(0,7)$, verified by checking all forty-nine products.

\subsection{Three ladder operators}

Pair the seven generators, with one left over, and complexify. Define three operators
\begin{equation}
  \alpha_k = \tfrac12\left(L_{2k-1} + i\,L_{2k}\right), \qquad k = 1, 2, 3 .
\end{equation}
From the Clifford relations we compute their algebra and find the canonical fermionic anticommutators,
\begin{equation}
  \{\alpha_j, \alpha_k^\dagger\} = \delta_{jk}, \qquad \{\alpha_j, \alpha_k\} = 0 .
\end{equation}
These are creation and annihilation operators for three fermionic modes. They build a Fock space, a vacuum $\Omega$
annihilated by every $\alpha_k$, and states made by acting with the $\alpha_k^\dagger$. With three modes the Fock
space has $2^3 = 8$ states. This eight is the generation, and it equals the eight of $\mathrm{Spin}(8)$, the same eight
that was forced at the pinch-point.

\subsection{Color and charge, read off}

Group the eight states by how many creation operators built them, the eigenvalue of the number operator
$N = \sum_k \alpha_k^\dagger \alpha_k$. We compute the spectrum and the multiplicities directly,
\begin{equation}
  N \in \{0, 1, 2, 3\}, \qquad \text{multiplicities } 1, 3, 3, 1 .
\end{equation}
The three modes are the three colors, and the multiplicities are exactly the color representations of one generation.
The single $N=0$ state is a color singlet. The three $N=1$ states are a color triplet. The three $N=2$ states are a
color anti-triplet. The single $N=3$ state is a color singlet. This is the color content of one generation, with no
extra states and none missing.

The electric charge is the number operator divided by three,
\begin{equation}
  Q = \tfrac{N}{3} \in \left\{ 0, \ \tfrac13, \ \tfrac23, \ 1 \right\}.
\end{equation}
We read off the particles. The $N=0$ singlet with $Q=0$ is the neutrino. The $N=1$ color triplet with $Q=\tfrac13$ is
the down quark in three colors. The $N=2$ color anti-triplet with $Q=\tfrac23$ is the up quark in three colors. The
$N=3$ singlet with $Q=1$ is the electron. The fractional charges, the long-standing puzzle of why quarks carry thirds,
are here a triviality. Charge is one third of a count of three color modes, so it comes in thirds. The denominator
three is the number of colors, which is the dimension of the imaginary quaternions inside the octonions.

\begin{result}[One generation of fermions]
The minimal spinor of $C\ell(0,7)$, built from octonion left-multiplications, is a three-mode fermionic Fock space
whose eight states are exactly one generation of the Standard Model, a lepton doublet and a color triplet of quarks,
with color multiplicities $1, 3, 3, 1$ and electric charges $0, \tfrac13, \tfrac23, 1$. The fractional charges are
forced, being a count of color modes divided by three.
\end{result}

We verify every number above by direct matrix computation, the Clifford relations, the ladder anticommutators, the
number-operator spectrum, the multiplicities, and the charges, all reported by the experiment suite.

\subsection{Chirality and the two half-spinors}

The eight-state Fock space is a spinor, and it is worth seeing how it carries handedness, the fact that the weak force
acts on left-handed matter and not right-handed. The spinor of $\mathrm{Spin}(8)$ is reducible into two halves, the two
inequivalent half-spinors $8_s$ and $8_c$, distinguished by a parity. In the weight description of $D_4$, the spinor
weights are the half-integer vectors
\begin{equation}
  \left( \pm\tfrac12, \pm\tfrac12, \pm\tfrac12, \pm\tfrac12 \right), \qquad 2^4 = 16 \ \text{in all},
\end{equation}
and these sixteen split by the sign of the product of their entries. The eight with an even number of minus signs are
$8_s$, the eight with an odd number are $8_c$. In the Fock-space picture this is the split by the parity of the number
operator, the even-occupancy states against the odd-occupancy states. Acting with one ladder operator flips the parity,
carrying $8_s$ to $8_c$, so the two halves are genuine mirror images and not the same representation. This is the
left and right of the world, and they are different objects, which is why a force can couple to one and not the other.

The three eight-dimensional representations close the geometry once more. The vector weights $8_v$ are the eight axis
vectors $(\pm 1, 0, 0, 0)$, and the two half-spinors are the sixteen half-integer vectors, so together
\begin{equation}
  8_v \ \oplus\ 8_s \ \oplus\ 8_c \ = \ 8 + 8 + 8 \ = \ 24,
\end{equation}
the twenty-four vertices of the dual 24-cell. The three faces of triality, the vector and the two spinors, are once
again a 24-cell. The handedness of matter and the lattice of spacetime are the same object seen two ways.

The full generation, with both chiralities and the antiparticles, doubles this through the complexification $\mathbb{C}
\otimes \mathbb{O}$, which carries the $C\ell(6)$ on which the gauge group acts in the next section.
